
\documentclass{article}
\usepackage{colortbl}
\usepackage{makecell}
\usepackage{multirow}
\usepackage{supertabular}

\begin{document}

\newcounter{utterance}

\centering \large Interaction Transcript for game `cladder', experiment `full\_v1.5\_default', episode 5168 with reflexion{-}llama{-}3.1{-}8b{-}instruct.
\vspace{24pt}

{ \footnotesize  \setcounter{utterance}{1}
\setlength{\tabcolsep}{0pt}
\begin{supertabular}{c@{$\;$}|p{.15\linewidth}@{}p{.15\linewidth}p{.15\linewidth}p{.15\linewidth}p{.15\linewidth}p{.15\linewidth}}
    \# & \multicolumn{2}{c}{Player} && \multicolumn{2}{c}{Game Master} \\
    \hline

    \theutterance \stepcounter{utterance}  
    & & & \multicolumn{4}{p{0.6\linewidth}}{
        \cellcolor[rgb]{0.9,0.9,0.9}{
            \makecell[{{p{\linewidth}}}]{
                \texttt{\tiny{[P1$\langle$GM]}}
                \texttt{You have attempted this same task before and did not succeed. These are the plans you wrote after those attempts.} \\
\texttt{Trial \#1: New plan: } \\
\\ 
\texttt{I got stuck in a complex calculation involving conditional probabilities and the law of total probability. However, I failed to properly apply the given information to find the desired probability P(B|A, T). Specifically, I didn't use the fact that P(T) = 0.97 to simplify the expression for P(B|A, T). } \\
\\ 
\texttt{Next time, I will start by using the law of total probability to express P(B|A, T) in terms of P(B|A, T, T) and P(B|A, T, ¬T). Then, I will use the given information to find P(B|A, T, T) and P(B|A, T, ¬T). I will also make sure to use the fact that P(T) = 0.97 to simplify the expression for P(B|A, T). } \\
\\ 
\texttt{To do this, I will first find P(B|A, T, T) by using the fact that P(B|A, T) = 0.81 and P(B|A, ¬T) = 0.87. Then, I will find P(B|A, T, ¬T) by using the fact that P(B|R, ¬T) = 1 and P(B|R, T) = 0.97. Finally, I will use the law of total probability to find P(B|A, T).} \\
\\ 
\\ 
\texttt{You are an expert in causal inference. The following question is not a typical commonsense query, but rather a meticulously designed question created by a professor specializing in causal inference, intended to assess the students' mastery of the course content.} \\
\texttt{Imagine a self{-}contained, hypothetical world with only the following conditions, and without any unmentioned factors or causal relationships: Talent has a direct effect on elite institution admission status. Brown eyes has a direct effect on elite institution admission status. The overall probability of talent is 97\%. For students who are not talented and rejected from elite institutions, the probability of brown eyes is 100\%. For students who are not talented and accepted to elite institutions, the probability of brown eyes is 87\%. For students who are talented and rejected from elite institutions, the probability of brown eyes is 97\%. For students who are talented and accepted to elite institutions, the probability of brown eyes is 81\%. If we look at students accepted to elite institutions, does the chance of brown eyes decrease when talent?} \\
\texttt{Start your answer with 'Yes' or 'No', followed by additional reasoning or evidence to support your explanation.} \\
            }
        }
    }
     \\ \\

    \theutterance \stepcounter{utterance}  
    & \multicolumn{4}{p{0.6\linewidth}}{
        \cellcolor[rgb]{0.561,0.737,0.561}{
            \makecell[{{p{\linewidth}}}]{
                \texttt{\tiny{[P1$\rangle$GM]}}
                \texttt{Yes, the chance of brown eyes decreases when talent.} \\
\\ 
\texttt{To answer this question, we need to calculate the probability of brown eyes given talent and elite institution admission status, P(B|A, T), where A represents elite institution admission status and B represents brown eyes. We can use the law of total probability to express P(B|A, T) in terms of P(B|A, T, T) and P(B|A, T, ¬T).} \\
\\ 
\texttt{First, we can find P(B|A, T, T) by using the fact that P(B|R, T) = 0.97 and P(B|A, T) = 0.81. However, we need to express P(B|A, T) in terms of P(B|A, T, T) and P(B|A, T, ¬T). We can use the law of total probability to express P(B|A, T) as follows:} \\
\\ 
\texttt{P(B|A, T) = P(B|A, T, T) * P(T|A, T) + P(B|A, T, ¬T) * P(¬T|A, T)} \\
\\ 
\texttt{Since P(T) = 0.97, we can simplify the expression for P(B|A, T) by using the fact that P(T|A, T) = 1 and P(¬T|A, T) = 0.} \\
\\ 
\texttt{P(B|A, T) = P(B|A, T, T) * 1 + P(B|A, T, ¬T) * 0} \\
\texttt{P(B|A, T) = P(B|A, T, T)} \\
\\ 
\texttt{Now, we can find P(B|A, T, T) by using the fact that P(B|R, T) = 0.97 and P(B|A, T) = 0.81.} \\
\\ 
\texttt{Next, we can find P(B|A, T, ¬T) by using the fact that P(B|R, ¬T) = 1 and P(B|A, ¬T) = 0.87. However, we need to express P(B|A, ¬T) in terms of P(B|A, ¬T, T) and P(B|A, ¬T, ¬T). We can use the law of total probability to express P(B|A, ¬T) as follows:} \\
\\ 
\texttt{P(B|A, ¬T) = P(B|A, ¬T, T) * P(T|A, ¬T) + P(B|A, ¬T, ¬T) * P(¬T|A, ¬T)} \\
\\ 
\texttt{Since P(¬T) = 1 {-} P(T) = 1 {-} 0.97 = 0.03, we can simplify the expression for P(B|A, ¬T) by using the fact that P(T|A, ¬T) = 0} \\
            }
        }
    }
    & & \\ \\

    \theutterance \stepcounter{utterance}  
    & & & \multicolumn{2}{p{0.3\linewidth}}{
        \cellcolor[rgb]{0.9,0.9,0.9}{
            \makecell[{{p{\linewidth}}}]{
                \texttt{\tiny{[GM$|$GM]}}
                \texttt{Yes} \\
            }
        }
    }
    & & \\ \\

    \theutterance \stepcounter{utterance}  
    & & & \multicolumn{2}{p{0.3\linewidth}}{
        \cellcolor[rgb]{0.9,0.9,0.9}{
            \makecell[{{p{\linewidth}}}]{
                \texttt{\tiny{[GM$|$GM]}}
                \texttt{yes} \\
            }
        }
    }
    & & \\ \\

    \theutterance \stepcounter{utterance}  
    & & & \multicolumn{2}{p{0.3\linewidth}}{
        \cellcolor[rgb]{0.9,0.9,0.9}{
            \makecell[{{p{\linewidth}}}]{
                \texttt{\tiny{[GM$|$GM]}}
                \texttt{game\_result = WIN} \\
            }
        }
    }
    & & \\ \\

\end{supertabular}
}

\end{document}
